% !TeX root = ../navier-stokes-manuscript.tex
The three-dimensional Navier--Stokes regularity problem asks whether a smooth solution can lose its smoothness in finite time.
The velocity \(u\) and pressure \(p\) must satisfy
\begin{equation}\label{eq:NavierStokesSystem}
  \begin{gathered}
    \partial_tu+(u\cdot\nabla)u+\nabla p=\nu\Delta u,
    \qquad
    \nabla\cdot u=0
    \quad\text{on }\R^3\times(0,\infty),
    \\
    u(\cdot,0)=u_0
    \quad\text{on }\R^3.
  \end{gathered}
\end{equation}
Here \(u\) is the \(\R^3\)-valued velocity, \(p\) is the scalar pressure, and \(\nu>0\) is the fixed kinematic viscosity.

Fefferman's whole-space formulation assumes a smooth, divergence-free initial velocity whose derivatives decay faster than every algebraic rate \parencite{Fefferman2000}.
We denote this class by \(\Ssigma\):
\begin{equation}\label{eq:FeffermanDatumClass}
  \Ssigma
  :=
  \left\{
  a\in C^\infty(\R^3;\R^3):
  \begin{array}{l}
    \nabla\cdot a=0,\\[0.3em]
    \forall\alpha\in\Nzero^3\;
    \forall K\in\Nzero\;
    \exists C_{\alpha,K}<\infty\;
    \forall x\in\R^3,\\[0.3em]
    |\partial_x^\alpha a(x)|
    \leq C_{\alpha,K}(1+|x|)^{-K}
  \end{array}
  \right\}.
\end{equation}
The constant \(C_{\alpha,K}\) may depend on \(a\), the multi-index \(\alpha\), and the decay order \(K\), while remaining independent of \(x\).
Choosing \(K>3/2\) for the derivatives up to any fixed order gives
\[
  H^m_\sigma(\R^3)
  :=\{v\in H^m(\R^3;\R^3):\nabla\cdot v=0\},
  \qquad
  \Ssigma\subset H^m_\sigma(\R^3)
  \quad(m\in\Nzero).
\]
In particular, \(u_0\) begins smooth, divergence-free, rapidly decreasing, and with finite kinetic energy.
Any later loss of regularity has to be generated by the Navier--Stokes evolution.

Fix \(u_0\in\Ssigma\) and \(\nu>0\).
The equation determines \(p\) up to an additive function of time.
Proposition~\ref{prop:LocalMaximalHistory} fixes this freedom by requiring \(p(\cdot,t)\in H^3(\R^3)\) and gives a unique maximal classical solution \((u,p)\) on \(\R^3\times[0,T_*)\), where \(T_*\in(0,\infty]\).
For every \(T<T_*\), the pair is smooth on \(\R^3\times[0,T]\) and solves \eqref{eq:NavierStokesSystem} with the same datum \(u_0\).
A continuation through a finite \(T_*\) would extend this solution to \([0,T_*+\varepsilon)\) for some \(\varepsilon>0\) and agree with it throughout \([0,T_*)\).
The continuation alternative in Proposition~\ref{prop:LocalMaximalHistory} states that
\[
  T_*<\infty
  \quad\Longrightarrow\quad
  \limsup_{t\uparrow T_*}\norm{u(t)}_{H^3}=\infty.
\]
A finite endpoint is accompanied by unbounded \(H^3\) velocity along the terminal approach.

The whole-space assertion asks that this maximal solution exist for all time and satisfy
\begin{equation}\label{eq:WholeSpaceRegularityTarget}
  u,p\in C^\infty\!\bigl(\R^3\times[0,\infty)\bigr),
  \qquad
  \sup_{t\geq0}\int_{\R^3}|u(x,t)|^2\,dx<\infty.
\end{equation}
For every \(0\leq t<T_*\), the classical energy identity gives \(\norm{u(t)}_2\leq\norm{u_0}_2\), so the energy in \eqref{eq:WholeSpaceRegularityTarget} is controlled throughout the maximal lifespan.
Local theory supplies smoothness on each compact subinterval of \([0,T_*)\).
The regularity problem is the assertion that this lifespan has no finite endpoint:
\[
  T_*=\infty.
\]
